%!TEX root = ../thesis.tex
% One MCTS simulation for \cref{sec:ez-mcts}, traced through the three phases.
% Worked example: |A| = 2 and gamma = 1 for legibility. The tree's Q values span
% exactly [0,1], so the tree-wide normalization of \cref{eq:ez-qnorm} is the
% identity here and the displayed values are both Q and \bar{Q}.
%
% Layout: the three phases are stacked as full-width rows rather than set side by
% side, and every tree grows LEFT TO RIGHT. Three 0.32\textwidth columns left only
% ~4.6cm per panel, which forced \tiny labels; one full-width row per phase gives
% ~8.7cm of drawing width and lets the statistics be set in \scriptsize. The
% subcaption sits beside its tree, not underneath it, so the rows stay short.
% Layout rule, to keep the panels collision-free: the simulation path runs
% straight right along y=0 with siblings on the row below; statistics of path
% nodes are therefore always set ABOVE, statistics of siblings always BELOW, and
% edge labels sit on the edge itself with a white background.
% Every panel opens with the same invisible \path at y=0, which pads all three
% tikzpictures to the same horizontal extent so s^0, s^1, s^2 line up in a column
% across the rows. It deliberately does not fix the vertical extent: (b) and (c)
% carry no statistics under their sibling nodes and would otherwise reserve the
% empty space that only (a) needs.
\begin{figure}[htbp]
  \centering
  % The subcaptions sit in a 0.36\textwidth column, too narrow to justify without
  % gaping word spaces -- set them ragged right for this figure only.
  \captionsetup[sub]{justification=raggedright, singlelinecheck=false}
  \tikzset{
    >={Stealth[round]},
    nd/.style      = {draw, circle, fill=black!8, minimum size=8mm,
                      inner sep=0pt, font=\small},
    ndpath/.style  = {nd, fill=blue!15, very thick},
    ndnew/.style   = {nd, fill=blue!15, very thick, draw=blue!60!black},
    ndopen/.style  = {draw, circle, densely dashed, black!50, minimum size=8mm,
                      inner sep=0pt, font=\small},
    net/.style     = {draw, rounded corners=2pt, fill=white, font=\scriptsize,
                      inner sep=2pt},
    ed/.style      = {->, black!55},
    edsel/.style   = {->, very thick, blue!60!black},
    edup/.style    = {->, very thick, densely dashed, black!75},
    pr/.style      = {font=\scriptsize, inner sep=1pt, black!55, fill=white},
    prsel/.style   = {font=\scriptsize, inner sep=1pt, fill=white},
    st/.style      = {font=\scriptsize, inner sep=1pt, align=center},
  }

  % ===== (a) Selection =====================================================
  \begin{subfigure}[b]{\textwidth}
    \begin{minipage}[c]{0.60\textwidth}
      \centering
      \begin{tikzpicture}
        \path (-1.1,0) (7.5,0);
        \node[ndpath] (s0) at (0,0)      {$s^0$};
        \node[ndpath] (s1) at (2.5,0)    {$s^1$};
        \node[ndopen] (nw) at (5.0,0)    {};
        \node[nd]     (A)  at (2.5,-1.5) {};
        \node[nd]     (B)  at (5.0,-1.5) {};

        \draw[edsel] (s0) -- node[prsel] {$p{=}0.75$} (s1);
        \draw[edsel] (s1) -- node[prsel] {$p{=}0.6$}  (nw);
        \draw[ed]    (s0) -- node[pr,pos=0.55] {$p{=}0.25$} (A);
        \draw[ed]    (s1) -- node[pr,pos=0.55] {$p{=}0.4$}  (B);

        \node[st,above=1.5mm of s0] {$N{=}6$};
        \node[st,above=1.5mm of s1] {$\bar{Q}{=}0.5$,\ $N{=}2$\\pUCT${}=\mathbf{1.27}$};
        \node[st,above=1.5mm of nw] {pUCT${}=\mathbf{1.06}$};
        \node[st,below=1.5mm of A]  {$\bar{Q}{=}1.0$,\ $N{=}4$\\pUCT${}=1.15$};
        \node[st,below=1.5mm of B]  {$\bar{Q}{=}0.0$,\ $N{=}2$\\pUCT${}=0.24$};
      \end{tikzpicture}
    \end{minipage}\hfill
    \begin{minipage}[c]{0.36\textwidth}
      \caption{Nodes are selected by pUCT (\cref{eq:ez-puct}).
        Despite action~1 having twice the action value of action~0,
        action~0 is chosen because of the higher prior and lower visit count.}
      \label{fig:ez-mcts-selection}
    \end{minipage}
  \end{subfigure}

  \par\vspace{3.5ex}
  % ===== (b) Expansion and evaluation ======================================
  \begin{subfigure}[b]{\textwidth}
    \begin{minipage}[c]{0.60\textwidth}
      \centering
      \begin{tikzpicture}
        \path (-1.1,0) (7.5,0);
        \node[ndpath] (s0) at (0,0)      {$s^0$};
        \node[ndpath] (s1) at (2.5,0)    {$s^1$};
        \node[ndnew]  (s2) at (5.0,0)    {$s^2$};
        \node[nd]     (A)  at (2.5,-1.5) {};
        \node[nd]     (B)  at (5.0,-1.5) {};
        \node[ndopen] (c0) at (7.0, 0.9) {};
        \node[ndopen] (c1) at (7.0,-0.9) {};

        \draw[edsel] (s0) -- (s1);
        \draw[ed]    (s0) -- (A);
        \draw[ed]    (s1) -- (B);
        \draw[edsel] (s1) -- node[net] {$g$} (s2);
        \draw[ed]    (s2) -- node[pr,pos=0.6] {$0.35$} (c0);
        \draw[ed]    (s2) -- node[pr,pos=0.6] {$0.65$} (c1);

        \node[net] (V) at (5.0,1.35) {$v$};
        \draw[ed] (s2) -- (V);
        \node[st,right=1.5mm of V] {$v^2{=}0.6$};
      \end{tikzpicture}
    \end{minipage}\hfill
    \begin{minipage}[c]{0.36\textwidth}
      % The \todo that used to sit here moved next to %!TEX root = ../thesis.tex
% One MCTS simulation for \cref{sec:ez-mcts}, traced through the three phases.
% Worked example: |A| = 2 and gamma = 1 for legibility. The tree's Q values span
% exactly [0,1], so the tree-wide normalization of \cref{eq:ez-qnorm} is the
% identity here and the displayed values are both Q and \bar{Q}.
%
% Layout: the three phases are stacked as full-width rows rather than set side by
% side, and every tree grows LEFT TO RIGHT. Three 0.32\textwidth columns left only
% ~4.6cm per panel, which forced \tiny labels; one full-width row per phase gives
% ~8.7cm of drawing width and lets the statistics be set in \scriptsize. The
% subcaption sits beside its tree, not underneath it, so the rows stay short.
% Layout rule, to keep the panels collision-free: the simulation path runs
% straight right along y=0 with siblings on the row below; statistics of path
% nodes are therefore always set ABOVE, statistics of siblings always BELOW, and
% edge labels sit on the edge itself with a white background.
% Every panel opens with the same invisible \path at y=0, which pads all three
% tikzpictures to the same horizontal extent so s^0, s^1, s^2 line up in a column
% across the rows. It deliberately does not fix the vertical extent: (b) and (c)
% carry no statistics under their sibling nodes and would otherwise reserve the
% empty space that only (a) needs.
\begin{figure}[htbp]
  \centering
  % The subcaptions sit in a 0.36\textwidth column, too narrow to justify without
  % gaping word spaces -- set them ragged right for this figure only.
  \captionsetup[sub]{justification=raggedright, singlelinecheck=false}
  \tikzset{
    >={Stealth[round]},
    nd/.style      = {draw, circle, fill=black!8, minimum size=8mm,
                      inner sep=0pt, font=\small},
    ndpath/.style  = {nd, fill=blue!15, very thick},
    ndnew/.style   = {nd, fill=blue!15, very thick, draw=blue!60!black},
    ndopen/.style  = {draw, circle, densely dashed, black!50, minimum size=8mm,
                      inner sep=0pt, font=\small},
    net/.style     = {draw, rounded corners=2pt, fill=white, font=\scriptsize,
                      inner sep=2pt},
    ed/.style      = {->, black!55},
    edsel/.style   = {->, very thick, blue!60!black},
    edup/.style    = {->, very thick, densely dashed, black!75},
    pr/.style      = {font=\scriptsize, inner sep=1pt, black!55, fill=white},
    prsel/.style   = {font=\scriptsize, inner sep=1pt, fill=white},
    st/.style      = {font=\scriptsize, inner sep=1pt, align=center},
  }

  % ===== (a) Selection =====================================================
  \begin{subfigure}[b]{\textwidth}
    \begin{minipage}[c]{0.60\textwidth}
      \centering
      \begin{tikzpicture}
        \path (-1.1,0) (7.5,0);
        \node[ndpath] (s0) at (0,0)      {$s^0$};
        \node[ndpath] (s1) at (2.5,0)    {$s^1$};
        \node[ndopen] (nw) at (5.0,0)    {};
        \node[nd]     (A)  at (2.5,-1.5) {};
        \node[nd]     (B)  at (5.0,-1.5) {};

        \draw[edsel] (s0) -- node[prsel] {$p{=}0.75$} (s1);
        \draw[edsel] (s1) -- node[prsel] {$p{=}0.6$}  (nw);
        \draw[ed]    (s0) -- node[pr,pos=0.55] {$p{=}0.25$} (A);
        \draw[ed]    (s1) -- node[pr,pos=0.55] {$p{=}0.4$}  (B);

        \node[st,above=1.5mm of s0] {$N{=}6$};
        \node[st,above=1.5mm of s1] {$\bar{Q}{=}0.5$,\ $N{=}2$\\pUCT${}=\mathbf{1.27}$};
        \node[st,above=1.5mm of nw] {pUCT${}=\mathbf{1.06}$};
        \node[st,below=1.5mm of A]  {$\bar{Q}{=}1.0$,\ $N{=}4$\\pUCT${}=1.15$};
        \node[st,below=1.5mm of B]  {$\bar{Q}{=}0.0$,\ $N{=}2$\\pUCT${}=0.24$};
      \end{tikzpicture}
    \end{minipage}\hfill
    \begin{minipage}[c]{0.36\textwidth}
      \caption{Nodes are selected by pUCT (\cref{eq:ez-puct}).
        Despite action~1 having twice the action value of action~0,
        action~0 is chosen because of the higher prior and lower visit count.}
      \label{fig:ez-mcts-selection}
    \end{minipage}
  \end{subfigure}

  \par\vspace{3.5ex}
  % ===== (b) Expansion and evaluation ======================================
  \begin{subfigure}[b]{\textwidth}
    \begin{minipage}[c]{0.60\textwidth}
      \centering
      \begin{tikzpicture}
        \path (-1.1,0) (7.5,0);
        \node[ndpath] (s0) at (0,0)      {$s^0$};
        \node[ndpath] (s1) at (2.5,0)    {$s^1$};
        \node[ndnew]  (s2) at (5.0,0)    {$s^2$};
        \node[nd]     (A)  at (2.5,-1.5) {};
        \node[nd]     (B)  at (5.0,-1.5) {};
        \node[ndopen] (c0) at (7.0, 0.9) {};
        \node[ndopen] (c1) at (7.0,-0.9) {};

        \draw[edsel] (s0) -- (s1);
        \draw[ed]    (s0) -- (A);
        \draw[ed]    (s1) -- (B);
        \draw[edsel] (s1) -- node[net] {$g$} (s2);
        \draw[ed]    (s2) -- node[pr,pos=0.6] {$0.35$} (c0);
        \draw[ed]    (s2) -- node[pr,pos=0.6] {$0.65$} (c1);

        \node[net] (V) at (5.0,1.35) {$v$};
        \draw[ed] (s2) -- (V);
        \node[st,right=1.5mm of V] {$v^2{=}0.6$};
      \end{tikzpicture}
    \end{minipage}\hfill
    \begin{minipage}[c]{0.36\textwidth}
      % The \todo that used to sit here moved next to %!TEX root = ../thesis.tex
% One MCTS simulation for \cref{sec:ez-mcts}, traced through the three phases.
% Worked example: |A| = 2 and gamma = 1 for legibility. The tree's Q values span
% exactly [0,1], so the tree-wide normalization of \cref{eq:ez-qnorm} is the
% identity here and the displayed values are both Q and \bar{Q}.
%
% Layout: the three phases are stacked as full-width rows rather than set side by
% side, and every tree grows LEFT TO RIGHT. Three 0.32\textwidth columns left only
% ~4.6cm per panel, which forced \tiny labels; one full-width row per phase gives
% ~8.7cm of drawing width and lets the statistics be set in \scriptsize. The
% subcaption sits beside its tree, not underneath it, so the rows stay short.
% Layout rule, to keep the panels collision-free: the simulation path runs
% straight right along y=0 with siblings on the row below; statistics of path
% nodes are therefore always set ABOVE, statistics of siblings always BELOW, and
% edge labels sit on the edge itself with a white background.
% Every panel opens with the same invisible \path at y=0, which pads all three
% tikzpictures to the same horizontal extent so s^0, s^1, s^2 line up in a column
% across the rows. It deliberately does not fix the vertical extent: (b) and (c)
% carry no statistics under their sibling nodes and would otherwise reserve the
% empty space that only (a) needs.
\begin{figure}[htbp]
  \centering
  % The subcaptions sit in a 0.36\textwidth column, too narrow to justify without
  % gaping word spaces -- set them ragged right for this figure only.
  \captionsetup[sub]{justification=raggedright, singlelinecheck=false}
  \tikzset{
    >={Stealth[round]},
    nd/.style      = {draw, circle, fill=black!8, minimum size=8mm,
                      inner sep=0pt, font=\small},
    ndpath/.style  = {nd, fill=blue!15, very thick},
    ndnew/.style   = {nd, fill=blue!15, very thick, draw=blue!60!black},
    ndopen/.style  = {draw, circle, densely dashed, black!50, minimum size=8mm,
                      inner sep=0pt, font=\small},
    net/.style     = {draw, rounded corners=2pt, fill=white, font=\scriptsize,
                      inner sep=2pt},
    ed/.style      = {->, black!55},
    edsel/.style   = {->, very thick, blue!60!black},
    edup/.style    = {->, very thick, densely dashed, black!75},
    pr/.style      = {font=\scriptsize, inner sep=1pt, black!55, fill=white},
    prsel/.style   = {font=\scriptsize, inner sep=1pt, fill=white},
    st/.style      = {font=\scriptsize, inner sep=1pt, align=center},
  }

  % ===== (a) Selection =====================================================
  \begin{subfigure}[b]{\textwidth}
    \begin{minipage}[c]{0.60\textwidth}
      \centering
      \begin{tikzpicture}
        \path (-1.1,0) (7.5,0);
        \node[ndpath] (s0) at (0,0)      {$s^0$};
        \node[ndpath] (s1) at (2.5,0)    {$s^1$};
        \node[ndopen] (nw) at (5.0,0)    {};
        \node[nd]     (A)  at (2.5,-1.5) {};
        \node[nd]     (B)  at (5.0,-1.5) {};

        \draw[edsel] (s0) -- node[prsel] {$p{=}0.75$} (s1);
        \draw[edsel] (s1) -- node[prsel] {$p{=}0.6$}  (nw);
        \draw[ed]    (s0) -- node[pr,pos=0.55] {$p{=}0.25$} (A);
        \draw[ed]    (s1) -- node[pr,pos=0.55] {$p{=}0.4$}  (B);

        \node[st,above=1.5mm of s0] {$N{=}6$};
        \node[st,above=1.5mm of s1] {$\bar{Q}{=}0.5$,\ $N{=}2$\\pUCT${}=\mathbf{1.27}$};
        \node[st,above=1.5mm of nw] {pUCT${}=\mathbf{1.06}$};
        \node[st,below=1.5mm of A]  {$\bar{Q}{=}1.0$,\ $N{=}4$\\pUCT${}=1.15$};
        \node[st,below=1.5mm of B]  {$\bar{Q}{=}0.0$,\ $N{=}2$\\pUCT${}=0.24$};
      \end{tikzpicture}
    \end{minipage}\hfill
    \begin{minipage}[c]{0.36\textwidth}
      \caption{Nodes are selected by pUCT (\cref{eq:ez-puct}).
        Despite action~1 having twice the action value of action~0,
        action~0 is chosen because of the higher prior and lower visit count.}
      \label{fig:ez-mcts-selection}
    \end{minipage}
  \end{subfigure}

  \par\vspace{3.5ex}
  % ===== (b) Expansion and evaluation ======================================
  \begin{subfigure}[b]{\textwidth}
    \begin{minipage}[c]{0.60\textwidth}
      \centering
      \begin{tikzpicture}
        \path (-1.1,0) (7.5,0);
        \node[ndpath] (s0) at (0,0)      {$s^0$};
        \node[ndpath] (s1) at (2.5,0)    {$s^1$};
        \node[ndnew]  (s2) at (5.0,0)    {$s^2$};
        \node[nd]     (A)  at (2.5,-1.5) {};
        \node[nd]     (B)  at (5.0,-1.5) {};
        \node[ndopen] (c0) at (7.0, 0.9) {};
        \node[ndopen] (c1) at (7.0,-0.9) {};

        \draw[edsel] (s0) -- (s1);
        \draw[ed]    (s0) -- (A);
        \draw[ed]    (s1) -- (B);
        \draw[edsel] (s1) -- node[net] {$g$} (s2);
        \draw[ed]    (s2) -- node[pr,pos=0.6] {$0.35$} (c0);
        \draw[ed]    (s2) -- node[pr,pos=0.6] {$0.65$} (c1);

        \node[net] (V) at (5.0,1.35) {$v$};
        \draw[ed] (s2) -- (V);
        \node[st,right=1.5mm of V] {$v^2{=}0.6$};
      \end{tikzpicture}
    \end{minipage}\hfill
    \begin{minipage}[c]{0.36\textwidth}
      % The \todo that used to sit here moved next to %!TEX root = ../thesis.tex
% One MCTS simulation for \cref{sec:ez-mcts}, traced through the three phases.
% Worked example: |A| = 2 and gamma = 1 for legibility. The tree's Q values span
% exactly [0,1], so the tree-wide normalization of \cref{eq:ez-qnorm} is the
% identity here and the displayed values are both Q and \bar{Q}.
%
% Layout: the three phases are stacked as full-width rows rather than set side by
% side, and every tree grows LEFT TO RIGHT. Three 0.32\textwidth columns left only
% ~4.6cm per panel, which forced \tiny labels; one full-width row per phase gives
% ~8.7cm of drawing width and lets the statistics be set in \scriptsize. The
% subcaption sits beside its tree, not underneath it, so the rows stay short.
% Layout rule, to keep the panels collision-free: the simulation path runs
% straight right along y=0 with siblings on the row below; statistics of path
% nodes are therefore always set ABOVE, statistics of siblings always BELOW, and
% edge labels sit on the edge itself with a white background.
% Every panel opens with the same invisible \path at y=0, which pads all three
% tikzpictures to the same horizontal extent so s^0, s^1, s^2 line up in a column
% across the rows. It deliberately does not fix the vertical extent: (b) and (c)
% carry no statistics under their sibling nodes and would otherwise reserve the
% empty space that only (a) needs.
\begin{figure}[htbp]
  \centering
  % The subcaptions sit in a 0.36\textwidth column, too narrow to justify without
  % gaping word spaces -- set them ragged right for this figure only.
  \captionsetup[sub]{justification=raggedright, singlelinecheck=false}
  \tikzset{
    >={Stealth[round]},
    nd/.style      = {draw, circle, fill=black!8, minimum size=8mm,
                      inner sep=0pt, font=\small},
    ndpath/.style  = {nd, fill=blue!15, very thick},
    ndnew/.style   = {nd, fill=blue!15, very thick, draw=blue!60!black},
    ndopen/.style  = {draw, circle, densely dashed, black!50, minimum size=8mm,
                      inner sep=0pt, font=\small},
    net/.style     = {draw, rounded corners=2pt, fill=white, font=\scriptsize,
                      inner sep=2pt},
    ed/.style      = {->, black!55},
    edsel/.style   = {->, very thick, blue!60!black},
    edup/.style    = {->, very thick, densely dashed, black!75},
    pr/.style      = {font=\scriptsize, inner sep=1pt, black!55, fill=white},
    prsel/.style   = {font=\scriptsize, inner sep=1pt, fill=white},
    st/.style      = {font=\scriptsize, inner sep=1pt, align=center},
  }

  % ===== (a) Selection =====================================================
  \begin{subfigure}[b]{\textwidth}
    \begin{minipage}[c]{0.60\textwidth}
      \centering
      \begin{tikzpicture}
        \path (-1.1,0) (7.5,0);
        \node[ndpath] (s0) at (0,0)      {$s^0$};
        \node[ndpath] (s1) at (2.5,0)    {$s^1$};
        \node[ndopen] (nw) at (5.0,0)    {};
        \node[nd]     (A)  at (2.5,-1.5) {};
        \node[nd]     (B)  at (5.0,-1.5) {};

        \draw[edsel] (s0) -- node[prsel] {$p{=}0.75$} (s1);
        \draw[edsel] (s1) -- node[prsel] {$p{=}0.6$}  (nw);
        \draw[ed]    (s0) -- node[pr,pos=0.55] {$p{=}0.25$} (A);
        \draw[ed]    (s1) -- node[pr,pos=0.55] {$p{=}0.4$}  (B);

        \node[st,above=1.5mm of s0] {$N{=}6$};
        \node[st,above=1.5mm of s1] {$\bar{Q}{=}0.5$,\ $N{=}2$\\pUCT${}=\mathbf{1.27}$};
        \node[st,above=1.5mm of nw] {pUCT${}=\mathbf{1.06}$};
        \node[st,below=1.5mm of A]  {$\bar{Q}{=}1.0$,\ $N{=}4$\\pUCT${}=1.15$};
        \node[st,below=1.5mm of B]  {$\bar{Q}{=}0.0$,\ $N{=}2$\\pUCT${}=0.24$};
      \end{tikzpicture}
    \end{minipage}\hfill
    \begin{minipage}[c]{0.36\textwidth}
      \caption{Nodes are selected by pUCT (\cref{eq:ez-puct}).
        Despite action~1 having twice the action value of action~0,
        action~0 is chosen because of the higher prior and lower visit count.}
      \label{fig:ez-mcts-selection}
    \end{minipage}
  \end{subfigure}

  \par\vspace{3.5ex}
  % ===== (b) Expansion and evaluation ======================================
  \begin{subfigure}[b]{\textwidth}
    \begin{minipage}[c]{0.60\textwidth}
      \centering
      \begin{tikzpicture}
        \path (-1.1,0) (7.5,0);
        \node[ndpath] (s0) at (0,0)      {$s^0$};
        \node[ndpath] (s1) at (2.5,0)    {$s^1$};
        \node[ndnew]  (s2) at (5.0,0)    {$s^2$};
        \node[nd]     (A)  at (2.5,-1.5) {};
        \node[nd]     (B)  at (5.0,-1.5) {};
        \node[ndopen] (c0) at (7.0, 0.9) {};
        \node[ndopen] (c1) at (7.0,-0.9) {};

        \draw[edsel] (s0) -- (s1);
        \draw[ed]    (s0) -- (A);
        \draw[ed]    (s1) -- (B);
        \draw[edsel] (s1) -- node[net] {$g$} (s2);
        \draw[ed]    (s2) -- node[pr,pos=0.6] {$0.35$} (c0);
        \draw[ed]    (s2) -- node[pr,pos=0.6] {$0.65$} (c1);

        \node[net] (V) at (5.0,1.35) {$v$};
        \draw[ed] (s2) -- (V);
        \node[st,right=1.5mm of V] {$v^2{=}0.6$};
      \end{tikzpicture}
    \end{minipage}\hfill
    \begin{minipage}[c]{0.36\textwidth}
      % The \todo that used to sit here moved next to \input{figures/ez-mcts-simulation}
      % in sections/efficientzero.tex: \todo uses \marginpar, and a \marginpar inside
      % a subfigure (a minipage) makes LaTeX abort with ``Float(s) lost''.
      \caption{The selected node is expanded. The latent state is rolled out once
        via $g$ and value statistics are initialized.}
      \label{fig:ez-mcts-expansion}
    \end{minipage}
  \end{subfigure}

  \par\vspace{3.5ex}
  % ===== (c) Backup ========================================================
  \begin{subfigure}[b]{\textwidth}
    \begin{minipage}[c]{0.60\textwidth}
      \centering
      \begin{tikzpicture}
        \path (-1.1,0) (7.5,0);
        \node[ndpath] (s0) at (0,0)      {$s^0$};
        \node[ndpath] (s1) at (2.5,0)    {$s^1$};
        \node[ndnew]  (s2) at (5.0,0)    {$s^2$};
        \node[nd]     (A)  at (2.5,-1.5) {};
        \node[nd]     (B)  at (5.0,-1.5) {};

        \draw[ed] (s0) -- (A);
        \draw[ed] (s1) -- (B);
        \draw[edup] (s1) -- node[prsel] {$r{=}0.2$} (s0);
        \draw[edup] (s2) -- node[prsel] {$r{=}0.4$} (s1);

        \node[st,above=1.5mm of s0] {$G^0{=}1.2$\\$N{:}\,6{\to}7$\\$Q_\Sigma{:}\,3.6{\to}4.8$};
        \node[st,above=1.5mm of s1] {$G^1{=}1.0$\\$N{:}\,2{\to}3$\\$Q_\Sigma{:}\,0.6{\to}1.6$};
        \node[st,above=1.5mm of s2] {$G^2{=}v^2{=}0.6$\\$N{:}\,0{\to}1$\\$Q_\Sigma{:}\,0{\to}0.6$};
      \end{tikzpicture}
    \end{minipage}\hfill
    \begin{minipage}[c]{0.36\textwidth}
      \caption{The leaf estimate is backed up through the path. Each node accumulates
        its return and the visit count is incremented.}
      \label{fig:ez-mcts-backup}
    \end{minipage}
  \end{subfigure}

  \caption[One MCTS simulation]{A single iteration of Monte-Carlo
    tree search, where $\lvert\mathcal{A}\rvert=2$ and $\gamma=1$.
    Search starts from a true root latent $s^0$ and is then continued by imagined rollouts
    to imagined latents $s^1$ and $s^2$.
    $Q_\Sigma$ represents the sum of all backed-up Q values, $N$ represents the visit count,
    $\bar{Q}=\frac{Q_\Sigma}{N}$ the average Q value of the node, and $\text{pUCT}$ the actual score
    according to the pUCT rule~(\cref{eq:ez-puct}).
    Colored nodes and edges indicate the search path while the edges are labeled with the prior $p$.}
  \label{fig:ez-mcts-simulation}
\end{figure}

      % in sections/efficientzero.tex: \todo uses \marginpar, and a \marginpar inside
      % a subfigure (a minipage) makes LaTeX abort with ``Float(s) lost''.
      \caption{The selected node is expanded. The latent state is rolled out once
        via $g$ and value statistics are initialized.}
      \label{fig:ez-mcts-expansion}
    \end{minipage}
  \end{subfigure}

  \par\vspace{3.5ex}
  % ===== (c) Backup ========================================================
  \begin{subfigure}[b]{\textwidth}
    \begin{minipage}[c]{0.60\textwidth}
      \centering
      \begin{tikzpicture}
        \path (-1.1,0) (7.5,0);
        \node[ndpath] (s0) at (0,0)      {$s^0$};
        \node[ndpath] (s1) at (2.5,0)    {$s^1$};
        \node[ndnew]  (s2) at (5.0,0)    {$s^2$};
        \node[nd]     (A)  at (2.5,-1.5) {};
        \node[nd]     (B)  at (5.0,-1.5) {};

        \draw[ed] (s0) -- (A);
        \draw[ed] (s1) -- (B);
        \draw[edup] (s1) -- node[prsel] {$r{=}0.2$} (s0);
        \draw[edup] (s2) -- node[prsel] {$r{=}0.4$} (s1);

        \node[st,above=1.5mm of s0] {$G^0{=}1.2$\\$N{:}\,6{\to}7$\\$Q_\Sigma{:}\,3.6{\to}4.8$};
        \node[st,above=1.5mm of s1] {$G^1{=}1.0$\\$N{:}\,2{\to}3$\\$Q_\Sigma{:}\,0.6{\to}1.6$};
        \node[st,above=1.5mm of s2] {$G^2{=}v^2{=}0.6$\\$N{:}\,0{\to}1$\\$Q_\Sigma{:}\,0{\to}0.6$};
      \end{tikzpicture}
    \end{minipage}\hfill
    \begin{minipage}[c]{0.36\textwidth}
      \caption{The leaf estimate is backed up through the path. Each node accumulates
        its return and the visit count is incremented.}
      \label{fig:ez-mcts-backup}
    \end{minipage}
  \end{subfigure}

  \caption[One MCTS simulation]{A single iteration of Monte-Carlo
    tree search, where $\lvert\mathcal{A}\rvert=2$ and $\gamma=1$.
    Search starts from a true root latent $s^0$ and is then continued by imagined rollouts
    to imagined latents $s^1$ and $s^2$.
    $Q_\Sigma$ represents the sum of all backed-up Q values, $N$ represents the visit count,
    $\bar{Q}=\frac{Q_\Sigma}{N}$ the average Q value of the node, and $\text{pUCT}$ the actual score
    according to the pUCT rule~(\cref{eq:ez-puct}).
    Colored nodes and edges indicate the search path while the edges are labeled with the prior $p$.}
  \label{fig:ez-mcts-simulation}
\end{figure}

      % in sections/efficientzero.tex: \todo uses \marginpar, and a \marginpar inside
      % a subfigure (a minipage) makes LaTeX abort with ``Float(s) lost''.
      \caption{The selected node is expanded. The latent state is rolled out once
        via $g$ and value statistics are initialized.}
      \label{fig:ez-mcts-expansion}
    \end{minipage}
  \end{subfigure}

  \par\vspace{3.5ex}
  % ===== (c) Backup ========================================================
  \begin{subfigure}[b]{\textwidth}
    \begin{minipage}[c]{0.60\textwidth}
      \centering
      \begin{tikzpicture}
        \path (-1.1,0) (7.5,0);
        \node[ndpath] (s0) at (0,0)      {$s^0$};
        \node[ndpath] (s1) at (2.5,0)    {$s^1$};
        \node[ndnew]  (s2) at (5.0,0)    {$s^2$};
        \node[nd]     (A)  at (2.5,-1.5) {};
        \node[nd]     (B)  at (5.0,-1.5) {};

        \draw[ed] (s0) -- (A);
        \draw[ed] (s1) -- (B);
        \draw[edup] (s1) -- node[prsel] {$r{=}0.2$} (s0);
        \draw[edup] (s2) -- node[prsel] {$r{=}0.4$} (s1);

        \node[st,above=1.5mm of s0] {$G^0{=}1.2$\\$N{:}\,6{\to}7$\\$Q_\Sigma{:}\,3.6{\to}4.8$};
        \node[st,above=1.5mm of s1] {$G^1{=}1.0$\\$N{:}\,2{\to}3$\\$Q_\Sigma{:}\,0.6{\to}1.6$};
        \node[st,above=1.5mm of s2] {$G^2{=}v^2{=}0.6$\\$N{:}\,0{\to}1$\\$Q_\Sigma{:}\,0{\to}0.6$};
      \end{tikzpicture}
    \end{minipage}\hfill
    \begin{minipage}[c]{0.36\textwidth}
      \caption{The leaf estimate is backed up through the path. Each node accumulates
        its return and the visit count is incremented.}
      \label{fig:ez-mcts-backup}
    \end{minipage}
  \end{subfigure}

  \caption[One MCTS simulation]{A single iteration of Monte-Carlo
    tree search, where $\lvert\mathcal{A}\rvert=2$ and $\gamma=1$.
    Search starts from a true root latent $s^0$ and is then continued by imagined rollouts
    to imagined latents $s^1$ and $s^2$.
    $Q_\Sigma$ represents the sum of all backed-up Q values, $N$ represents the visit count,
    $\bar{Q}=\frac{Q_\Sigma}{N}$ the average Q value of the node, and $\text{pUCT}$ the actual score
    according to the pUCT rule~(\cref{eq:ez-puct}).
    Colored nodes and edges indicate the search path while the edges are labeled with the prior $p$.}
  \label{fig:ez-mcts-simulation}
\end{figure}

      % in sections/efficientzero.tex: \todo uses \marginpar, and a \marginpar inside
      % a subfigure (a minipage) makes LaTeX abort with ``Float(s) lost''.
      \caption{The selected node is expanded. The latent state is rolled out once
        via $g$ and value statistics are initialized.}
      \label{fig:ez-mcts-expansion}
    \end{minipage}
  \end{subfigure}

  \par\vspace{3.5ex}
  % ===== (c) Backup ========================================================
  \begin{subfigure}[b]{\textwidth}
    \begin{minipage}[c]{0.60\textwidth}
      \centering
      \begin{tikzpicture}
        \path (-1.1,0) (7.5,0);
        \node[ndpath] (s0) at (0,0)      {$s^0$};
        \node[ndpath] (s1) at (2.5,0)    {$s^1$};
        \node[ndnew]  (s2) at (5.0,0)    {$s^2$};
        \node[nd]     (A)  at (2.5,-1.5) {};
        \node[nd]     (B)  at (5.0,-1.5) {};

        \draw[ed] (s0) -- (A);
        \draw[ed] (s1) -- (B);
        \draw[edup] (s1) -- node[prsel] {$r{=}0.2$} (s0);
        \draw[edup] (s2) -- node[prsel] {$r{=}0.4$} (s1);

        \node[st,above=1.5mm of s0] {$G^0{=}1.2$\\$N{:}\,6{\to}7$\\$Q_\Sigma{:}\,3.6{\to}4.8$};
        \node[st,above=1.5mm of s1] {$G^1{=}1.0$\\$N{:}\,2{\to}3$\\$Q_\Sigma{:}\,0.6{\to}1.6$};
        \node[st,above=1.5mm of s2] {$G^2{=}v^2{=}0.6$\\$N{:}\,0{\to}1$\\$Q_\Sigma{:}\,0{\to}0.6$};
      \end{tikzpicture}
    \end{minipage}\hfill
    \begin{minipage}[c]{0.36\textwidth}
      \caption{The leaf estimate is backed up through the path. Each node accumulates
        its return and the visit count is incremented.}
      \label{fig:ez-mcts-backup}
    \end{minipage}
  \end{subfigure}

  \caption[One MCTS simulation]{A single iteration of Monte-Carlo
    tree search, where $\lvert\mathcal{A}\rvert=2$ and $\gamma=1$.
    Search starts from a true root latent $s^0$ and is then continued by imagined rollouts
    to imagined latents $s^1$ and $s^2$.
    $Q_\Sigma$ represents the sum of all backed-up Q values, $N$ represents the visit count,
    $\bar{Q}=\frac{Q_\Sigma}{N}$ the average Q value of the node, and $\text{pUCT}$ the actual score
    according to the pUCT rule~(\cref{eq:ez-puct}).
    Colored nodes and edges indicate the search path while the edges are labeled with the prior $p$.}
  \label{fig:ez-mcts-simulation}
\end{figure}
